\chapter{Trabajo Futuro}
\label{chap:trabajo_futuro}

\lettrine{L}{a} hipótesis central del trabajo intenta demostrar que las tres familias de técnicas son
complementarias (Capítulo~\ref{chap:conclusions}). Cada familia presenta ventajas frente a las otras; 
la vía más prometedora no es perfeccionar ninguna por separado, sino combinarlas.
Este capítulo recoge aproximaciones futuras que se podrían implementar tras los descubrimientos
de este trabajo.

\paragraph{Inicialización del refuerzo por imitación}
\label{sec:tf-il-rl}
En lugar de entrenar las técnicas de refuerzo desde cero, se puede utilizar la política de imitación
ya entrenada (Capítulo~\ref{chap:iter2}). La red del agente de refuerzo se inicializa 
con los pesos del agente de \gls{il} y el refuerzo 
se encarga únicamente de refinar ese comportamiento y de cubrir lo que la
imitación no aprendió, realizando así un proceso de \gls{fine-tuning}. 

Esta inicialización como arranque del \gls{rl} ataca de frente el problema observado en el
Capítulo~\ref{chap:resultados}: el refuerzo desde cero invierte la mayor parte del cómputo
en aprender lo que la imitación ya tiene aprendido. En el proceso de investigación no se 
encontró literatura que aplicara esta hibridación específica de \gls{rl} con \gls{il} 
en Minecraft.

\paragraph{El planificador como guía de la exploración}
\label{sec:tf-htn-rl}
La segunda hibridación usa el \gls{htn} (Capítulo~\ref{chap:iter1}) como 
\textbf{tutor de exploración} en el entrenamiento. La idea es similar a la de \gls{dagger}, 
pero aplicada al refuerzo: el agente itera sobre el entorno y, 
cuando se desvíe o no sepa qué hacer, el planificador le indica la acción
que él habría tomado, alimentando así su aprendizaje con trayectorias útiles.
Combinada con un algoritmo \emph{off-policy} como \gls{sac}, 
esta guía puede integrarse directamente en el
\gls{replay-buffer}, mezclando transiciones del tutor con las del propio agente.

\paragraph{Ampliación de la cobertura del experto y subredes}
\label{sec:tf-cobertura}
El mayor problema presente en la imitación es la limitación del \gls{dataset}. Al no tener ejemplos de
exploración, el agente no puede aprender a buscar. Como solución se plantea el uso de \textit{subdatasets};
con ellos se podría dividir la tarea en subobjetivos y entrenar subredes especializadas en cada uno de ellos,
las cuales se irían intercambiando entre sí según la situación. Por ejemplo, una subred especializada en \emph{explorar}
se mantendría activa hasta que el agente encontrara un árbol, 
momento en el que se activaría la subred de \emph{aproximación} al árbol, dando paso a la subred de \emph{talado} y, por
último, a la subred de \emph{recogida}.

De esta forma, la recolección del \gls{dataset} se podría paralelizar al poder lanzar varios agentes con tareas más cortas,
abarcando así todas las situaciones posibles que el agente de imitación 
podría encontrarse en la tarea de \gls{woodcutting}.

\paragraph{Escalado de la tarea}
\label{sec:tf-escalado}
Todas las técnicas se evaluaron sobre la tarea de \gls{woodcutting}, una de las tareas
más simples de la jerarquía de progresión de Minecraft (Capítulo~\ref{chap:introduccion}). El siguiente
paso es ascender por esa jerarquía (fabricar una mesa de trabajo, un pico de
madera, recolectar piedra) hasta tareas de horizonte mucho más largo,
como la obtención del pico de hierro o de diamantes. 

\paragraph{Introducción de otras técnicas de aprendizaje}
\label{sec:tf-otras-tecnicas}
Si bien las familias de técnicas estudiadas abarcan un amplio espectro de enfoques, 
existen otras técnicas que también presentan un fuerte potencial. Como se menciona en el Capítulo~\ref{chap:sota}, el proyecto
MindCraft~\cite{mindcraft2025} intenta completar todo el juego (una tarea con un horizonte extremadamente
largo) utilizando únicamente modelos de lenguaje \gls{llm}, sin ningún tipo de intervención humana más
allá de la inicialización. Además, incluye varias aproximaciones como \gls{mas}, en la que múltiples agentes
trabajan en conjunto para resolver la tarea.

Dada la capacidad de los \gls{llm} para generar código y sus altas ventanas de contexto, la posibilidad
de completar buena parte de la jerarquía de progresión de Minecraft con estas técnicas es 
un área de investigación muy prometedora.